\begin{topic}[id=wio_terminal_tinyml,
             difficulty={Mittel},
             type={Praxisnaher Systementwurf},
             prerequisites={Grundlagen Mikrocontroller; Basisverständnis Sensorik und TinyML},
             practice={empfohlen},
             owner={Dozententeam},
             updated={2026-04-09},
             tags={ Wio Terminal, Edge Impulse, Datensatz, Klassifikation, MCU, Visualisierung }]{Wio Terminal + TinyML}

\TopicDescription{Der Seeed Wio Terminal kombiniert Sensoren, Display und MCU in einem günstigen Lern-Setup. Du siehst den kompletten Weg von der Datenerhebung über Feature-Engineering bis zur Inferenz auf dem Gerät. Wichtig sind die typischen Hürden kleiner Hardware (Speicher, Stabilität, Rauschen) und wie man Ergebnisse sichtbar macht. Nach dem Thema kannst du Aufwand und Grenzen solcher Demos realistisch einschätzen und nächste Schritte (Datennachschärfung, Modellvereinfachung, Energiemessung) planen.}

\TopicLeitfragen{\item Welche Schritte im Edge-Impulse-Workflow sind kritisch?
\item Wie visualisiert man Ergebnisse sinnvoll auf dem Gerät?
\item Welche typischen Fehlerquellen lassen sich früh vermeiden?
}
\TopicScope{\item Kein umfassendes Hardware-Tuning.
\item Keine tiefe Audio-/Signalverarbeitung.
}
\TopicPractice{Gesten- oder Klopferkennung mit Anzeige der Entscheidung auf dem Display.}

\TopicKeywords{Wio Terminal, Edge Impulse, Datensatz, Klassifikation, MCU, Visualisierung}

\nocite{seeedWioTinyml,edgeImpulseWio}
\end{topic}
